\documentclass[11pt]{article}

\usepackage[a4paper,margin=2.4cm]{geometry}
\usepackage{amsmath,amssymb}
\usepackage{siunitx}
\usepackage{xcolor}
\usepackage[hidelinks]{hyperref}
\usepackage{enumitem}

\title{Simulating the PET measurement\\
\large A user guide to the \texttt{PTCryspMC.jl} detector simulation}
\author{}
\date{\today}

\begin{document}
\maketitle

\section{Introduction}
\label{sec:intro}

\paragraph{Purpose.}
The simulation produces, for a given PET scanner, the list of coincidences it would
record from the positron activity left by a proton field. This list is the input to a
separate analysis that measures how precisely the proton range can be recovered. That
analysis is independent of the simulation and is not described here.

\paragraph{Inputs.}
\begin{itemize}[nosep]
\item A \textbf{scenario}, produced upstream by \texttt{ptcryspg4} and read from
\texttt{ptcrysp-scenarios}: the set of annihilation points and, per isotope, the number
of decays measured (the budget) for a given dose and acquisition timing.
\item A \textbf{detector description}: the ring geometry (radius, axial length, crystal
thickness), the crystal material, and the energy, position and time resolutions.
\end{itemize}

\paragraph{Approach.}
A photon-only Monte Carlo. Each annihilation emits two \SI{511}{keV} photons
back-to-back; the simulation transports them through the phantom and the detector ring,
records where they interact, and pairs them into coincidences. Only the photons are
tracked: the electrons they free deposit their energy where they are made (sub-millimetre
range at \SI{511}{keV}), and scintillation light is not modelled. This reproduces the
detector response while staying fast enough to run the order of $10^{8}$ annihilations a
measurement contains, for each candidate scanner.

\section{Structure of the simulation}
\label{sec:structure}

\subsection{Input photons}
\label{sec:input}

A scenario gives the annihilation points and, per isotope, the measured count $N_j$. The
photon set is built by drawing $N_j$ annihilation points per isotope from the scenario,
with replacement: the points supply the spatial distribution of the activity, $N_j$
supplies the count. For a \SI{1}{Gy} in-room field $\sum_j N_j$ is of order $10^{8}$.

Each drawn point emits two \SI{511}{keV} photons in a random direction, back-to-back, with
an acollinearity of about $0.5^\circ$ FWHM. The positron range is already contained in the
annihilation point (set upstream), so the emission point needs no further treatment.

\subsection{Detector simulation}
\label{sec:detsim}

A photon is transported until it is absorbed or leaves the setup.

\paragraph{Interactions.}
The distance to the next interaction is sampled from the material's total cross section
(NIST XCOM data, tabulated per material). At the interaction point the process is chosen
in proportion to the partial cross sections. Two are relevant at \SI{511}{keV}:
\begin{itemize}[nosep]
\item Compton scattering: the photon transfers energy to an electron and continues with
reduced energy in a new direction, both sampled from the Klein--Nishina distribution; the
electron energy is deposited at the interaction point.
\item Photoelectric absorption: the photon is absorbed and its energy deposited at the
interaction point.
\end{itemize}
Pair production is below threshold at \SI{511}{keV}.

\paragraph{Geometry.}
The phantom is a cylinder. The detector is a ring of crystal, segmented into blocks around
the ring (in $\varphi$) and stacked into wheels along the axis ($z$) --- the units a real
scanner reads out. We model this as a structured grid on a cylindrical shell, not as a
thousand separate volumes: the block a point belongs to is fixed by simple arithmetic on
its $\varphi$ and $z$, and the distance to the next boundary is the nearest of the inner
and outer radius, the next $\varphi$ plane and the next $z$ plane --- all closed-form.
Stepping a photon through the segmented ring therefore costs the same as through a single
shell, whatever the number of blocks. The small dead gaps between blocks are skipped for
now (the crystal is treated as continuous and each deposit is tagged with its block); they
can be added later as thin non-interacting layers. Transport advances a photon by the
smaller of the distance to the next boundary and the distance to the next interaction.

Through the setup a photon crosses the phantom (where it may Compton-scatter --- the
patient-scatter background --- or pass through), an air gap, and the crystal, where it
interacts at some depth. The depth of the first crystal interaction is the depth of
interaction (DOI).

\paragraph{In the crystal: overspill and containment.}
Three things can happen in the crystal, and all of them come out of following the photon.
It can deposit its full energy in one block. It can Compton-scatter and have the scattered
photon cross into a neighbouring block before being absorbed --- an overspill, leaving
energy in two blocks. Or the scattered photon can leave the crystal altogether (out the
back, the axial end, or back toward the patient) before depositing the rest, so only part
of the \SI{511}{keV} is recorded. A photon may also Compton-scatter more than once within
the same block. The transport produces all of these by construction; the bookkeeping below
turns the deposits into hits and selects events.

\paragraph{Hits and detector response.}
The deposits are grouped by block, and those in one block are reduced to a hit: the energy
is their sum, the position is the \emph{first} interaction point, and the time is the
earliest. The first interaction --- not the energy-weighted centroid --- is the point a
line of response should pass through, and the quantity a reconstruction aims to recover
(in a monolithic crystal a CNN can separate one- from two-Compton events to find it). The
hit is then blurred by the detector resolution: position, including DOI, by $\sigma_{xyz}$;
time by $\sigma_t$; and energy by an energy-dependent width
$\mathrm{FWHM}(E)=a\sqrt{\SI{511}{keV}/E}$, with $a$ the FWHM at \SI{511}{keV} ---
\SI{5}{\percent} for pure CsI, \SI{7}{\percent} for CsI(Tl), \SI{10}{\percent} for
cryogenic BGO, the monolithic crystals considered here. (Pixelated detectors --- LYSO,
standard BGO --- report a fixed crystal rather than a continuous position, and are left for
later.) A photon that overspilled leaves a hit in more than one block; the full list of
interactions in a block is kept in the singles record, so the one- versus two-Compton
separation can be studied later. A hit is kept if its energy falls in a symmetric window
about \SI{511}{keV}, of half-width $\sim$2 FWHM and tunable; because the width scales with
the resolution, a sharper detector rejects more patient scatter. These resolutions and the
window are the detector parameters; changing them is how candidate scanners differ.

\paragraph{One annihilation at a time.}
The measured rate is modest and spread thin. For a \SI{1}{Gy} field the total is of order
$10^{8}$ decays, an average of $\sim$$10^{5}$/s over the acquisition; the rate is
front-loaded, peaking near $7\times10^{5}$/s at the start of the window and falling with
$^{15}$O's \SI{122}{s} half-life. Spread over the $\sim$$10^{3}$ crystal modules of the
ring, each module sees $\sim$$10^{2}$/s, so even with a slow ($\sim$\SI{1}{\micro s})
crystal the chance of two signals overlapping in one module is $\sim$$10^{-4}$. Events
therefore do not pile up, and the simulation processes annihilations independently, one
pair of photons at a time.

\paragraph{Same-annihilation coincidences.}
For each annihilation we look at the blocks it lit up. The clean case is exactly two
touched blocks, roughly opposite, one hit per photon, each passing the energy window: the
pair is recorded as a coincidence --- true if neither photon scattered in the phantom,
scatter if one did. Anything else fails the condition and the event yields nothing
(efficiency loss): a photon that missed the ring, one that escaped the crystal with too
little energy, or one that overspilled so a third block was touched. (Recovering overspill
by clustering adjacent blocks is a possible refinement.) Every detected hit is also written
to a singles list, with its position, energy, parent isotope and scatter flag.

\paragraph{Randoms.}
A random --- two photons from different annihilations that happen to coincide --- cannot
occur within a single-annihilation event, so randoms are added in a separate pass over
the singles list, with no further transport. Each single is given a time drawn from the
activity curve of its parent isotope (which reproduces the front-loading); the singles
are sorted in time, and any two from different annihilations that fall within the
coincidence window $\tau$, pass the energy window and form a valid line of response are
recorded as randoms. Their rate follows $2\tau S(t)^2$ with $S$ the total singles rate,
so it is largest at the start of the scan and falls as the square of the activity; for
$\tau\sim\SI{10}{ns}$ it is at most a few to ten percent of the trues.

\paragraph{Running once.}
The transport is the costly part and is run once, producing the same-annihilation
coincidences and the singles list. The randoms pass --- and any re-realisation of the
measurement the downstream analysis may need --- are cheap operations on the singles
list and require no further transport.

\subsection{Output}
\label{sec:output}

The output is the coincidence list, written as \texttt{coincidences\_<config>.csv}: per
accepted coincidence, the two three-dimensional hit positions, the two energies, the two
times, and the truth label (true, scatter or random). The singles list is an
intermediate, kept to build the randoms. The coincidence list is what the downstream
analysis reads, and the only thing it reads.

\section{Implementation}
\label{sec:impl}

The photon transport and the geometry already exist in LXeMC, a liquid-xenon detector
simulation that does the same kind of \SI{}{keV}-scale photon transport, and are the
starting point. The pieces to reuse:
\begin{itemize}[nosep]
\item \textbf{Geometry} (\texttt{geometry\_core.jl}): the solids \texttt{Cyl} and
\texttt{CylShell}, with \texttt{is\_inside} and the ray entry/exit distances
(\texttt{distance\_to\_entry}, \texttt{distance\_to\_exit}). The phantom is one
\texttt{Cyl}; the ring is one \texttt{CylShell} with a $(\varphi, z)$ block index and the
corresponding plane-crossing distances added on top --- the one piece of new geometry.
\item \textbf{Cross sections} (\texttt{nist\_data.jl}, \texttt{materials.jl}): the XCOM
loader and \texttt{sigma\_three(material, E)}, which returns the Compton, photoelectric and
pair cross sections at an energy. New material tables are required: water or soft tissue
for the phantom, and the crystal materials (CsI, LYSO, BGO).
\item \textbf{Transport} (\texttt{sampling.jl}, \texttt{tracking.jl}): the Compton and
photoelectric samplers and the photon step loop (\texttt{transport\_photon!},
\texttt{propagate\_gamma}). The electron transport and bremsstrahlung in
\texttt{tracking.jl} are not needed here --- only the photon path with local energy
deposition.
\end{itemize}

What must be written on top:
\begin{itemize}[nosep]
\item \textbf{Source injection}: read a scenario, draw annihilation points per isotope to
the budget, and emit the back-to-back \SI{511}{keV} pairs with acollinearity
(Sec.~\ref{sec:input}).
\item \textbf{Hit formation} --- group deposits by block, take the first-interaction point,
sum the energy, smear, apply the energy window --- and the same-annihilation selection of
two opposite blocks (Sec.~\ref{sec:detsim}).
\item \textbf{The randoms pass}: time-tag the singles from the per-isotope activity, then
pair cross-annihilation singles within $\tau$ (Sec.~\ref{sec:detsim}). It reads the
singles list and writes randoms into the coincidence list, with no transport.
\end{itemize}

The transport runs in Julia. The source injection, the hit formation and smearing, the
coincidence sorting and the randoms pass --- reading a scenario and writing the
coincidence list --- run in Python.

\end{document}
